% Build with: tectonic paper.tex
\documentclass[11pt,letterpaper]{article}
\usepackage[T1]{fontenc}
\usepackage{lmodern}
\usepackage[margin=1in,headheight=15pt]{geometry}
\usepackage{amsmath,amssymb,amsthm}
\usepackage{microtype,booktabs,enumitem,xcolor,fancyhdr,hyperref}
\usepackage[most]{tcolorbox}
\usepackage{array}
\usepackage{needspace}
\newcolumntype{P}[1]{>{\raggedright\arraybackslash}p{#1}}
\definecolor{ink}{HTML}{203B55}
\definecolor{accent}{HTML}{176C65}
\definecolor{boxbg}{HTML}{F2F6F8}
\hypersetup{colorlinks=true,linkcolor=ink,urlcolor=accent,citecolor=ink,
 pdftitle={Depth-Two Threshold Circuits for Andreev's Function: Sharper Bounds},
 pdfauthor={Christian Reitwiessner}}
\pagestyle{fancy}
\fancyhf{}
\fancyhead[L]{\small\color{ink}Threshold circuits for Andreev's function}
\fancyhead[R]{\small\thepage}
\renewcommand{\headrulewidth}{0.3pt}
\setlength{\parindent}{0pt}
\setlength{\parskip}{0.55em}
\setlength{\emergencystretch}{2em}
\setlist[itemize]{leftmargin=1.5em,itemsep=0.2em,topsep=0.2em}
\setlist[enumerate]{leftmargin=1.8em,itemsep=0.2em,topsep=0.2em}
\newtheorem{theorem}{Theorem}
\newtheorem{lemma}[theorem]{Lemma}
\newtheorem{corollary}[theorem]{Corollary}
\newtheorem{proposition}[theorem]{Proposition}
\newtheorem{fact}[theorem]{Fact}
\theoremstyle{remark}
\newtheorem*{remark}{Remark}
\newcommand{\bits}{\{0,1\}}
\newcommand{\ind}[1]{\bigl[\,#1\,\bigr]}
\newcommand{\sect}[1]{\section{\texorpdfstring{\color{ink}#1}{#1}}}
\newcommand{\subsect}[1]{\subsection*{\texorpdfstring{\color{ink}#1}{#1}}}
\newcommand{\smalltitle}[1]{\textbf{\color{ink}#1}\quad}
\newcommand{\THR}{\mathsf{THR}}
\newcommand{\cH}{\mathcal H}
\newtcolorbox{keybox}[1][]{breakable,colback=boxbg,colframe=ink!40,boxrule=0.4pt,
  arc=2pt,left=8pt,right=8pt,top=5pt,bottom=5pt,#1}
\begin{document}
\thispagestyle{empty}
{\color{accent}\small RESEARCH NOTE\par}
\vspace{0.4em}
{\raggedright\color{ink}\LARGE\bfseries Depth-Two Threshold Circuits for\\[0.15em]
Andreev's Function: Sharper Bounds\par}
\vspace{0.5em}
{\large Christian Reitwiessner\par}
{\small November 2026\par}
\vspace{0.8em}

\begin{keybox}
\textbf{Abstract.} Andreev's function on $n$ input bits reads a bit from a
table at an address given by the parities of blocks of input bits. For
depth-two threshold circuits we show:
\begin{itemize}
\item an exact circuit with $(5/16+o(1))n^2$ gates and polynomial integer
      weights (previously $O(n^3/\log n)$ by Kane and Williams~\cite{KW});
\item for agreement on a $\tfrac12+\varepsilon$ fraction of inputs, matching
      bounds of $\varepsilon^2n^{3/2}$ gates up to polylogarithmic factors,
      for $\varepsilon$ between about $\log n/\sqrt n$ and $\tfrac18$;
      the upper bound is $O(\varepsilon^2n^{3/2}\sqrt{\log(1/\varepsilon)})$,
      and constant error needs only $O(n^{3/2})$ gates;
\item worst-case lower bounds of $\Omega(n^{3/2}/\log^2 n)$ gates and
      $\Omega(n^{5/2}/\log^{5/2}n)$ wires (previously $\log^3 n$ and
      $\log^{7/2}n$ in the denominators).
\end{itemize}
In depth three we improve the top fan-in bound for Andreev's function. For
the depth-three function $B_n$ of Kane and Williams we outline improved
lower bounds, saving one logarithm, together with an average-case bound that
was implicit in their work, for $B_n$ with explicit parameters. The lower
bounds for Andreev's function, an earlier quadratic circuit, and earlier
average-case upper bounds have merged Lean formalizations. The two sharper
upper bounds proved here are not yet formalized; precise coverage and the
status of the depth-three results are listed in Section~\ref{sec:lean}. A preliminary note (Appendix~\ref{app:coherent})
shows that circuits whose gates have a single sign inside every address block
need $\Omega(n^2/\log^3n)$ gates.
\end{keybox}

%===================================================================
\sect{Setting and results}

\smalltitle{Andreev's function.}
Fix integers $k,m\ge1$ and put $N=2^k$. The input consists of
\begin{itemize}
\item a \emph{table} $x$ of $N$ bits $x_c$, one for each address
      $c\in\bits^k$, and
\item $k$ \emph{address blocks}; block $i$ consists of the bits
      $a_{i,1},\dots,a_{i,m}$. Write $a$ for all address bits and
      $|a_i|$ for the number of ones in block $i$.
\end{itemize}
The \emph{address} $b\in\bits^k$ has coordinates $b_i=|a_i|\bmod2$, and
\[
A(x,a)=x_b .
\]
The input length is $n=N+km$. In the \emph{balanced case} $km=N$, so $n=2N$
and $k=\log_2(n/2)$. Exact balance requires $k\mid2^k$; for a family at
every $k$ one may instead take $m=\lfloor N/k\rfloor$, giving $n=2N-O(k)$.
Unless stated otherwise, asymptotic statements refer to these balanced regimes.

\smalltitle{Threshold circuits.}
For a statement $P$, the bracket $\ind{P}$ is $1$ if $P$ holds and $0$
otherwise. A \emph{threshold gate} with inputs $z_j\in\bits$, real weights
$\omega_j$ and threshold $\tau$ computes $\ind{\sum_j\omega_jz_j\ge\tau}$. In a \emph{depth-two threshold
circuit} ($\THR\circ\THR$) some \emph{bottom} gates read input bits, and a
single \emph{top} gate reads the bottom gates and input bits. Its
\emph{size} is its total number of gates, top gate included; its
\emph{wires} are its nonzero weights. A circuit
\emph{$\varepsilon$-agrees} with $A$ if it outputs $A(x,a)$ on at least a
$\tfrac12+\varepsilon$ fraction of all inputs.

\smalltitle{Results and the state of the art.}
Stronger lower bounds against $\THR\circ\THR$ are known for functions in the
much larger class $\mathsf E^{\mathsf{NP}}$; they are proved by the
algorithmic method~\cite{Tam,ACW,CTW}. For functions in $\mathsf P$, which
includes Andreev's function and $B_n$, the best known gate lower bounds are of
order $n^{3/2}/\mathrm{polylog}\,n$~\cite{KW,KKL}, as also noted
in~\cite{CTW}. Our bounds improve the explicit polylogarithmic factors and
the dependence on $\varepsilon$.
\begin{center}
\small
\renewcommand{\arraystretch}{1.25}
\setlength{\tabcolsep}{5pt}
\begin{tabular}{@{}lllll@{}}
\toprule
Function & Bound & Best previously known & This paper & \\
\midrule
\multicolumn{5}{@{}l}{\emph{Functions in $\mathsf P$}}\\
Andreev, exact & gates, upper & $O(n^3/\log n)$ \cite{KW} & $(5/16+o(1))n^2$ & Thm.~\ref{thm:exact}\\
Andreev, exact & gates, lower & \begin{tabular}[t]{@{}l@{}}$\Omega(n^{3/2}/\log^3n)$ \cite{KW}\\ $n^{3/2}/(\log n)^{O(1)}$ \cite{KKL}$^\dagger$\end{tabular} & $\Omega(n^{3/2}/\log^2n)$ & Thm.~\ref{thm:lower}\\
Andreev, exact & wires, lower & $n^{5/2}/(\log n)^{O(1)}$ \cite{KKL}$^\dagger$ & $\Omega(n^{5/2}/\log^{5/2}n)$ & Thm.~\ref{thm:lower}\\
Andreev, $\varepsilon$-agree & gates, upper & none stated & $O(\varepsilon^2n^{3/2}\sqrt{\log(1/\varepsilon)})$ & Thm.~\ref{thm:avgupper}\\
Andreev, $\varepsilon$-agree & gates, lower & $\Omega(\varepsilon^3n^{3/2}/\log^3n)$ \cite{KW} & $\Omega(\varepsilon^2n^{3/2}/\log^2n)$ & Thm.~\ref{thm:lower}\\
Andreev, $\varepsilon$-agree & wires, lower & $\Omega(\varepsilon^3n^{5/2}/\log^{7/2}n)$ \cite{KW} & $\Omega(\varepsilon^2n^{5/2}/\log^{5/2}n)$ & Thm.~\ref{thm:lower}\\
Andreev, depth 3 & top fan-in & $\tilde\Omega(n^{1/2}/s^{1/3})$, implicit in~\cite{KW} & $\min\{\tilde\Omega(n^{1/2}),\tilde\Omega(n^{3/4}/s^{1/2})\}$ & Cor.~\ref{cor:fanin}\\
$B_n$, depth 3, exact & gates, lower & $\Omega(n^{3/2}/\log^3n)$ \cite{KW} & $\Omega(n^{3/2}/\log^2n)$ & Thm.~\ref{thm:B}$^*$\\
$B_n$, depth 3, exact & wires, lower & $\Omega(n^{5/2}/\log^{7/2}n)$ \cite{KW} & $\Omega(n^{5/2}/\log^{5/2}n)$ & Thm.~\ref{thm:B}$^*$\\
$B_n$, $\varepsilon$-agree & gates, lower & $\tilde\Omega(\varepsilon n^{3/2})$, implicit in~\cite{KW} & $\Omega(\varepsilon n^{3/2}/\log^2n)$ & Thm.~\ref{thm:B}$^*$\\
\midrule
\multicolumn{5}{@{}l}{\emph{Functions in $\mathsf{E}^{\mathsf{NP}}$}}\\
some $F$, exact & gates, lower & \begin{tabular}[t]{@{}l@{}}$n^{2-\delta}$ \cite{Tam,ACW}\\ $n^{5/2-\delta}$ \cite{CTW}\end{tabular} & -- & \\
some $F$, $\varepsilon$-agree & wires, lower & $n^3/\log^{12}n$ \cite{Nag} (unreviewed) & -- & \\
\bottomrule
\end{tabular}
\end{center}
Notes on the table:
\begin{itemize}
\item All rows concern $\THR\circ\THR$ circuits unless they say depth~3.
      The Andreev rows are asymptotic in the balanced case; its exact upper
      bound has the explicit finite version \eqref{eq:bestcount}.
\item The $\varepsilon$-agreement rows for Andreev's function hold for
      $3\sqrt2(\log_2n+1)/\sqrt n\le\varepsilon$, and the upper bound for
      $\varepsilon\le\tfrac18$. For smaller $\varepsilon$, a single
      majority gate over the table already agrees on a
      $\tfrac12+\Omega(1/\sqrt n)$ fraction.
\item In the top fan-in row, $s$ bounds the size of each depth-two
      subcircuit.
\item $^\dagger$For a variant of Andreev's function whose table is four times
      as long as the address part.
\item $^*$For $B_n$ with the parameters of Theorem~\ref{thm:B}; the
      $\varepsilon$-agreement row holds for $\varepsilon\ge2\cdot2^{-3p}$, about
      $n^{-3}$. Proved here as an outline with constants; the full proof is
      formalized (open pull request). Chen, Santhanam and
      Srinivasan~\cite[Thm.~9.3]{CSS} re-prove the Kane--Williams bound for
      $B_n$ with $n^{o(1)}$ losses; Kabanets, Kane and Lu~\cite[Thm.~1.2]{KKL}
      prove a related wire bound for a variant of $B_n$.
\item Because of the quadratic upper bound, no gate lower bound above
      $n^2$ holds for Andreev's function, so the $n^{5/2-\delta}$ bound
      of~\cite{CTW} cannot hold for it.
\end{itemize}
The remaining gap for Andreev's function is the worst case: between
$n^{3/2}/\log^2n$ and $(5/16+o(1))n^2$ gates.

\smalltitle{How to read this paper.}
Section~\ref{sec:tech} explains which techniques are new.
Sections~\ref{sec:upper}--\ref{sec:avgupper} (upper bounds) and
Section~\ref{sec:lower} (lower bounds) can be read independently;
Section~\ref{sec:depth3} uses Section~\ref{sec:lower}. Each lemma uses only
its own local variables plus the notation in the boxes, and each box says
where its notation is valid.

%===================================================================
\sect{Where the improvements come from}\label{sec:tech}

Each new result combines standard tools with one or two new observations.
This section is an overview; the terms are defined precisely later. Briefly:
a \emph{character} is a sign $(-1)^{S\cdot c}$ for a set $S$ of address
coordinates; a \emph{staircase} writes a function of an integer linear form as
a sum of threshold gates (Lemma~\ref{lem:ladder}); a \emph{restriction} fixes
all but one bit per address block, and a gate \emph{survives} if it still
depends on the remaining bits (Section~\ref{sec:lower}); for $B_n$, the
table is replaced by a \emph{seed} (Section~\ref{sec:depth3}).

\begin{center}
\small
\renewcommand{\arraystretch}{1.25}
\begin{tabular}{@{}P{0.39\textwidth}P{0.33\textwidth}P{0.2\textwidth}@{}}
\toprule
Technique & Status & Used for\\
\midrule
Fourier expansion of the table lookup & standard & Thms.~\ref{thm:exact},~\ref{thm:avgupper}\\
Telescoping a function of a linear form into threshold gates (``staircase'') & standard~\cite[Prop.~1]{HMPST}, \cite{HP} & Thms.~\ref{thm:exact},~\ref{thm:avgupper}\\
Cutting staircases to a window via concentration & \cite[Thm.~1.4]{KW}, for parity & Thm.~\ref{thm:avgupper}\\
Grouping by character and removing the parity factor (Lemmas~\ref{lem:lookup},~\ref{lem:parity}) & starting quadratic construction & Prop.~\ref{prop:basic}\\
Shared quadratic corrections and sign margins at exceptional addresses & refinement proved here & Thm.~\ref{thm:exact}\\
Independent address masks cancel cross-character error moments & mean-square refinement proved here & Thm.~\ref{thm:avgupper}\\
\textbf{Majority vote over a subcube} (sum over $S\subseteq\{1,\dots,k_0\}$) & \textbf{new in this context} & Thm.~\ref{thm:avgupper}\\
Random restrictions with Littlewood--Offord & \cite[Lemmas~1.1, 3.1]{KW} & Thm.~\ref{thm:lower}\\
Counting threshold functions over fixed features & \cite{RSO}, \cite[Thm.~2.3]{KW} & Thm.~\ref{thm:lower}\\
\textbf{Fixed-weight counting} (Lemma~\ref{lem:count}) & structural fact standard~\cite{KW,CSS}; \textbf{its use for counting new, to our knowledge} & one $\log$ in Thm.~\ref{thm:lower}\\
\textbf{Uniform-in-$q$ union bound with Cauchy--Schwarz} (Lemma~\ref{lem:avg}) & \textbf{new, to our knowledge}; \cite{KW,KKL,CSS} use Markov & $\varepsilon^2$ in Thm.~\ref{thm:lower}\\
Averaging over a family of circuits & \cite[Cor.~1.1]{KW} & Cor.~\ref{cor:fanin}\\
Nearly orthogonal vectors (Lemma~\ref{lem:gram}) & Johnson bound~\cite{Gur}; \cite[Lemma~6.1]{KW} & Thm.~\ref{thm:B}\\
\textbf{Fixing the circuit before the seed} & \textbf{new for $B_n$, to our knowledge}; a standard probabilistic move & Thm.~\ref{thm:B}\\
\bottomrule
\end{tabular}
\end{center}

\smalltitle{Upper bounds.}
Kane and Williams~\cite[Thm.~1.2]{KW} write $x_b$ as a threshold of about
$n^2/\log n$ parities and build each parity from its own depth-two
subcircuit of $O(n)$ gates. We group the parities that share the same
address character into one product of a linear form and a parity, then
remove the parity with an elementary identity (``removing the parity
factor'', Lemma~\ref{lem:parity}). This saves the factor $n/\log n$: a
generic simulation of a threshold of parities builds every parity separately,
whereas Lemma~\ref{lem:parity} lets all parities with the same address
character share one staircase. The sharper exact score collapses quadratic
corrections across characters and absorbs errors at a few special addresses
using sign margins. The average-case circuit keeps only the characters of
the first $r$ address bits, producing a majority vote over a subcube.
Independent fixed masks make the characters' truncation errors orthogonal
in expectation, avoiding a union bound over all characters.

\smalltitle{Lower bounds.}
We keep the restriction argument of Kane and Williams and change two counting
steps.
\begin{itemize}
\item \emph{Fixed-weight counting.} After the restriction, a surviving gate
      keeps its weights on the live bits; only its threshold changes. This
      fact is standard; it appears in the proof of~\cite[Lemma~1.1]{KW}, as
      equation~(9.1) of~\cite{CSS}, and in satisfiability algorithms for
      threshold circuits~\cite{CTW}. We use it for counting: a survivor costs
      $O(k+\log s)$ bits to describe, not the $\Theta(k^2)$ bits of an
      arbitrary threshold function, as in~\cite{KW,KKL,CSS}. This removes one
      logarithm. It works because the circuit is fixed before the random table
      is drawn; arguments that fix the hard input before the circuit, by
      counting~\cite{KW,CSS} or by incompressibility~\cite{KKL}, cannot do
      this.
\item \emph{Uniform union bound.} Kane and Williams bound the number of
      surviving gates with Markov's inequality, which costs a factor of
      $\varepsilon$. We bound the correlation for all survivor counts at
      once and average with Cauchy--Schwarz. This changes $\varepsilon^3$
      to $\varepsilon^2$.
\end{itemize}
For $B_n$, fixed-weight counting needs the circuit to be fixed before the
seed is chosen; Kane and Williams fix the seed first.

%===================================================================
\sect{Exact quadratic circuits}\label{sec:upper}

\begin{keybox}
\textbf{Notation for Sections~\ref{sec:upper} and~\ref{sec:avgupper}.}
For $S\subseteq\{1,\dots,k\}$ and $c\in\bits^k$ write $S\cdot c=\sum_{i\in S}c_i$, and
\[
X_S=\sum_{c\in\bits^k}(-1)^{S\cdot c+x_c},\qquad
t_S=\sum_{i\in S}|a_i| .
\]
$X_S$ depends only on the table, is even, and lies in $[-N,N]$.
$t_S$ is the number of ones in the blocks in $S$, and lies in $[0,|S|m]$.
\end{keybox}

\smalltitle{Plan.}
We first compute the integer score $N\cdot(-1)^{x_b}$ and test its sign.
Lemma~\ref{lem:lookup} writes it as a sum of products,
Lemma~\ref{lem:parity} removes the products, and Lemma~\ref{lem:ladder}
turns the result into threshold gates. We then refine this to a smaller score
with the same sign, proving the $5/16$ bound.

\begin{lemma}[Lookup as a sum of products]\label{lem:lookup}
$\displaystyle\sum_{S\subseteq\{1,\dots,k\}}X_S\,(-1)^{t_S}=N\,(-1)^{x_b}.$
\end{lemma}
\begin{proof}
Since $|a_i|\equiv b_i\pmod2$, we have $(-1)^{t_S}=(-1)^{S\cdot b}$.
Exchanging the two sums, the left side is
\[
\sum_c(-1)^{x_c}\prod_{i=1}^k\bigl(1+(-1)^{c_i+b_i}\bigr),
\]
and the product is $N$ for $c=b$ and $0$ otherwise.
\end{proof}

Let $\varphi(\theta)=(-1)^\theta\,\theta$ for integers $\theta$.

\begin{lemma}[Removing the parity factor]\label{lem:parity}
If $X$ is an even integer and $t$ is an integer, then
$X(-1)^t=\varphi(X+t)-\varphi(t)$.
\end{lemma}
\begin{proof}
$X+t$ has the parity of $t$, so $\varphi(X+t)=(-1)^t(X+t)$.
\end{proof}

\begin{lemma}[Threshold staircase]\label{lem:ladder}
Let $L$ be an integer that always lies in $\{\ell,\dots,h\}$, and let $g$
be any function on the integers. Then
\[
g(L)=g(\ell)+\sum_{\theta=\ell+1}^{h}\bigl(g(\theta)-g(\theta-1)\bigr)\ind{L\ge\theta}.
\]
\end{lemma}
\begin{proof}
The brackets equal $1$ exactly for $\theta\le L$, and the sum telescopes.
\end{proof}

When $L$ is a linear form in the input bits, each $\ind{L\ge\theta}$ is one
threshold gate. Both $t_S$ and $X_S+t_S$ are such forms, since
$(-1)^{x_c}=1-2x_c$:
\begin{equation}\label{eq:form}
X_S+t_S=\sum_c(-1)^{S\cdot c}-2\sum_c(-1)^{S\cdot c}x_c+\sum_{i\in S}\sum_{j=1}^m a_{i,j}.
\end{equation}

\begin{proposition}[A first quadratic construction]\label{prop:basic}
For all $k,m\ge1$, Andreev's function has an exact depth-two threshold
circuit with $2N^2+kmN$ bottom gates and one top gate. In the balanced case
this is $\tfrac34n^2+1$ gates. Bottom weights lie in $\{0,\pm2\}$ on table
bits and in $\{0,1\}$ on address bits, and
top weights are odd integers of absolute value at most $2n-1$.
\end{proposition}
\begin{proof}
By Lemmas~\ref{lem:lookup} and~\ref{lem:parity},
\[
N(-1)^{x_b}=\sum_S\bigl(\varphi(X_S+t_S)-\varphi(t_S)\bigr).
\]
Apply Lemma~\ref{lem:ladder} to every term. $X_S+t_S$ ranges over
$[-N,\,N+|S|m]$, which needs $2N+|S|m$ gates, and $t_S$ over $[0,|S|m]$,
which needs $|S|m$ gates. Each index $i$ lies in half of the $N$ sets $S$,
so $\sum_S|S|=kN/2$, and the number of bottom gates is
$\sum_S(2N+2|S|m)=2N^2+kmN$. The top gate weights the gate
$\ind{L\ge\theta}$ by $\pm(\varphi(\theta)-\varphi(\theta-1))=\pm(-1)^\theta(2\theta-1)$.
This is odd, and since $-N<\theta\le N+km=n$, its absolute value is at most
$2n-1$.
\end{proof}
The term $S=\varnothing$ is $X_\varnothing$, a linear form in the table that the
top gate could read directly; this would save $2N$ gates. We keep the uniform
construction.
Since the bottom weights are small integers, the circuit is even a threshold
of majority gates ($\THR\circ\mathsf{MAJ}$, by repeating inputs). For
comparison, Kane and Williams also give a $\mathsf{MOD}_3\circ\mathsf{MOD}_2$
circuit of size $O(n^2)$~\cite[Thm.~1.2]{KW}, in a different basis.

% Copyright (c) 2026 Christian Reitwiessner.
% Authors: Christian Reitwiessner
\subsect{A sharper score: the \texorpdfstring{$5/16$}{5/16} bound}

We can save more gates by asking only for the correct sign. The score below
recovers the selected bit with a fixed margin, while allowing a controlled
error at the zero address and at the $k$ addresses of Hamming weight one.
A shared correction handles the zero address; amplifying the corresponding
table entries absorbs the other errors.

\begin{theorem}\label{thm:exact}
Put $\alpha=\lfloor(m-1)/2\rfloor$. There is an exact depth-two threshold
circuit for $A$ with at most
\begin{equation}\label{eq:bestcount}
\frac{N(N-1)}2+\frac{kN(m+\alpha)}2
 +2(N-1)+2\alpha k+3km+\left\lceil\frac{km}{2}\right\rceil
\end{equation}
bottom gates and one top gate. Its weights are integers of magnitude
polynomial in $n$. At exact balance its total size is
\[
\frac5{16}n^2-\frac n4\log_2(n/2)+O(n).
\]
For $m=\lfloor N/k\rfloor$, its size is $(5/16+o(1))n^2$.
\end{theorem}

\begin{proof}
\emph{Forms used only in this proof.}
Write $s_i=|a_i|$, $\sigma=\sum_i s_i$, and let $e_i$ be the address with
only coordinate $i$ equal to one. Put
\[
z=\sum_{c\ne0}x_c,\quad E=\sum_i x_{e_i},\quad
w=2z+2\alpha E+\sigma,\quad R=2(N-1)+2\alpha k+km,\quad C=N(R+4).
\]
Thus $0\le w\le R$. For each nonempty $S\subseteq[k]$, define
\[
Y_S=\sum_{c:S\cdot c\text{ odd}}x_c,\quad
W_S=Y_S+\alpha\sum_{i\in S}x_{e_i},\quad
\kappa_S=-4(2\alpha+1)(k-2|S|),
\]
and let $r_S(u)=-8\lceil u/2\rceil^2+\kappa_S\lceil u/2\rceil$.
All of $W_S,t_S,w,s_i$ are affine forms in the original inputs.

\emph{The score.} Consider
\begin{align}\label{eq:bestscore}
P={}&2N+\sum_{S\ne\varnothing}r_S(2W_S+t_S)
 +\frac N2\left(w^2+\sum_i s_i^2\right)\notag\\
 &+4Nz-2N(\alpha^2+1)E-2\alpha N\sigma
 +2N(1+\alpha)\sum_i x_{e_i}s_i-Cx_0\ind{b=0}.
\end{align}
We first verify its sign, then realize its non-affine terms with bottom gates.
The following identities follow by expanding the squares and summing
characters; all sums over $S$ are over nonempty sets:
\begin{align*}
\sum_S4W_S(-1)^{t_S}
 &=2N\ind{b=0}z-2N\ind{b\ne0}x_b
   +2\alpha N\sum_i x_{e_i}\bigl(\ind{b=0}-\ind{b=e_i}\bigr),\\
8\sum_S(W_S^2+W_St_S)
 &=\frac N2(w^2-\sigma^2)+2Nz+(4\alpha+2\alpha^2)NE
   +2N(1+\alpha)\sum_i x_{e_i}s_i,\\
8\sum_S\lceil t_S/2\rceil^2
 &=\frac N2\left(\sigma^2+\sum_i s_i^2\right)
   +N(\sigma+1)\ind{b\ne0}+N\sum_i s_i\ind{b=e_i}.
\end{align*}
For example, the first is character orthogonality applied to $Y_S$ and to
$\ind{i\in S}$; the other two use the same identity and
$\lceil t/2\rceil=(t+\ind{t\text{ odd}})/2$.
Also
\[
4W(-1)^t=4W-8\lceil(2W+t)/2\rceil^2+8W^2+8Wt+8\lceil t/2\rceil^2.
\]
The coefficient shifts $\kappa_S$ supply the address correction because
\[
\sum_S\kappa_S\bigl(\lceil t_S/2\rceil-t_S/2\bigr)
 =N(2\alpha+1)\sum_i\ind{b=e_i}.
\]
Substitution in \eqref{eq:bestscore} now gives the simple closed form
\begin{equation}\label{eq:bestclosed}
P=\begin{cases}
N(w+2)-Cx_0,&b=0,\\
2N(1+\alpha)(1-x_{e_i})-Ns_i,&b=e_i,\\
N(1-2x_b),&\text{otherwise}.
\end{cases}
\end{equation}
At $b=0$, use $0\le w\le R$ and $C=N(R+4)$. At $b=e_i$, the odd
integer $s_i$ lies between $1$ and $2\alpha+1$. In every case $P\ge N$
when $x_b=0$, and $P\le-N$ when $x_b=1$. The top gate outputs $\ind{-P\ge0}$.

\emph{Raw gates and their count.}
The staircase for $r_S(2W_S+t_S)$ uses only odd thresholds
\[
\ind{2W_S+t_S\ge\theta},\qquad
1\le\theta\le N+2\alpha|S|+|S|m,
\]
with coefficient $-8\theta+\kappa_S$. The two square terms use $R+km$
gates by Lemma~\ref{lem:ladder}. For $B=n+2$, each product $x_{e_i}s_i$
uses $m$ gates $\ind{s_i+Bx_{e_i}\ge B+\ell}$, $1\le\ell\le m$.
Finally, the identity
\[
x_0\ind{b=0}=x_0-x_0\sigma+
 \frac4N\sum_{S\ne\varnothing}x_0\lfloor t_S/2\rfloor
\]
implements the last correction with $km$ gates
$\ind{\sigma+Bx_0\ge B+\ell}$ and
$\sum_S\lfloor |S|m/2\rfloor$ gates
$\ind{t_S+Bx_0\ge B+2j}$. Its identity follows from
$\sum_S\ind{t_S\text{ odd}}=(N/2)\ind{b\ne0}$.
The even-level $\sigma$ gates coincide with the $S=[k]$ gates, saving
$\lfloor km/2\rfloor$ gates. Everything else in $P$ is an affine top term.

Since $\sum_{S\ne\varnothing}|S|=kN/2$, the main-ridge and last-correction
counts together are
\[
\sum_{S\ne\varnothing}\left(
 \left\lceil\frac{N+2\alpha|S|+|S|m}{2}\right\rceil
 +\left\lfloor\frac{|S|m}{2}\right\rfloor\right)
 =\frac{N(N-1)}2+\frac{\alpha kN}2+\frac{kmN}2.
\]
Adding the remaining gates gives \eqref{eq:bestcount}. Vanishing coefficients
or further coincidences can only improve it. At exact balance $m$ is even,
so $\alpha=m/2-1$; substitution gives the stated asymptotic bound.
For floor balance, the rounding changes the leading expression by
$O(Nk)=o(n^2)$. Bottom coefficients and thresholds are $O(n)$; the displayed
expansion has top coefficients bounded by $O(n^3)$.
\end{proof}

The sharper construction remains sign-coherent: every bottom sees a block
through its sum with a nonnegative coefficient. Large or mixed-sign weights
are therefore not needed for this bound. Section~\ref{app:coherent} shows
that improving the quadratic exponent requires leaving that class.


%===================================================================
% Copyright (c) 2026 Christian Reitwiessner.
% Authors: Christian Reitwiessner
\sect{Agreeing on most inputs with fewer gates}\label{sec:avgupper}

The earlier windowed construction pays for keeping every character inside
its window at once. Here we bound the sum of the errors in mean square.
Independent address masks make cross-character error terms cancel, removing
the factor $\sqrt{\log n}$ for constant error.

\Needspace{12\baselineskip}
\begin{theorem}\label{thm:avgupper}
In the balanced case, the following depth-two threshold circuits exist:
\begin{enumerate}
\item For $0<\delta\le1/2$, agreement at least $1-\delta$ with
      $O\bigl(n^{3/2}\sqrt{\log(1/\delta)}\bigr)$ gates.
\item For $N^{-1/2}\le\varepsilon\le1/8$, agreement at least
      $1/2+\varepsilon$ with
      $O\bigl(\varepsilon^2n^{3/2}\sqrt{\log(1/\varepsilon)}\bigr)$ gates
      and $O\bigl(\varepsilon^2n^{5/2}\sqrt{\log(1/\varepsilon)}\bigr)$ wires.
\end{enumerate}
For $1/8<\varepsilon<1/2$, Part~1 applies with $\delta=1/2-\varepsilon$.
The proof chooses one fixed circuit for each size; it does not assert an
efficient uniform algorithm for choosing its masks.
\end{theorem}

\begin{proof}
\emph{A majority vote with fewer characters.}
Fix $0\le r\le k$ and put $q=2^r$. Character orthogonality gives
\[
P_*=2\sum_{S\subseteq[r]}X_S(-1)^{t_S}
 =2q\sum_{c:\,c_{[r]}=b_{[r]}}(-1)^{x_c}.
\]
Its sign is a majority vote of $L=N/q$ independent table bits, one of them
$x_b$. With ties outputting zero, its agreement is
$1/2+p(L)/2$, where $p(1)=1$ and
\[
p(L)=2^{-L}\binom{L}{L/2}\ge(2L)^{-1/2}\qquad(L\ge2).
\]
Indeed, conditioning on the selected bit leaves a sum of $L-1$ independent
signs; the excess agreement is half the probability that this sum equals
$-1$. This is the subcube-majority calculation used by the earlier circuit.

\emph{Masked, truncated staircases.}
Independently for each nonempty $S\subseteq[r]$, choose a uniform mask $z_S$
on its $M_S=|S|m$ active raw address coordinates. Let
\[
D_S=|a_S\oplus z_S|-\lfloor M_S/2\rfloor,\qquad
\eta_S=(-1)^{|z_S|+\lfloor M_S/2\rfloor}.
\]
With $\varphi(u)=(-1)^u u$, evenness of $X_S$ gives
\[
\eta_S\bigl(\varphi(X_S+D_S)+\varphi(X_S-D_S)\bigr)
 =2X_S(-1)^{t_S}.
\]
For an integer radius $R\ge1$, set
$f_R(u)=\varphi(u)\ind{|u|\le R}$, and replace the left side by its
$f_R$ version. Together with the free affine term $2X_\varnothing$, these
terms give a score $P_z$. The finite-difference expansion
\[
f_R(U)=\sum_{\ell=-R}^{R+1}
 \bigl(f_R(\ell)-f_R(\ell-1)\bigr)\ind{U\ge\ell}
\]
uses at most $2R+2$ gates. Each $X_S\pm D_S$ is an affine form in the raw
inputs, with table weights $\pm2$ and address weights $\pm1$.
Thus $P_z$ uses at most $4(q-1)(R+1)$ bottoms and one top gate.

\emph{Cancellation of cross terms.}
Let $e_S$ be the truncated contribution minus its exact contribution.
Conditional on $x,a$, the masks are independent, and
$D_S$ has a shifted binomial law independent of $a$. Moreover
$\eta_S=(-1)^{t_S}(-1)^{D_S}$. Therefore
\[
\mathbb E_{z_S}[e_S\mid x,a]=(-1)^{t_S}g_S(X_S)
\]
for a function $g_S$ of the table alone. For $S\ne T$, averaging the product
over masks and then uniform raw addresses gives
\[
\mathbb E_{z,x,a}[e_Se_T]
 =\mathbb E_x[g_S(X_S)g_T(X_T)]\,
   \mathbb E_a[(-1)^{t_S+t_T}]=0.
\]
No independence of Walsh coefficients is assumed. Consequently
\begin{equation}\label{eq:maskl2}
\mathbb E_{z,x,a}(P_z-P_*)^2=\sum_{S\ne\varnothing}\mathbb E e_S^2.
\end{equation}

\emph{Diagonal terms and one fixed circuit.}
For $U=X_S\pm D_S$, the mean has magnitude at most $1/2$ and the centered
variable is subgaussian with variance proxy $v_S=N+M_S/4$. Integrating its
Hoeffding tail yields
\[
\mathbb E\bigl[U^2\ind{|U|>R}\bigr]
 \le16v_S\exp\bigl(-R^2/(16v_S)\bigr).
\]
For example, integrate
$\Pr(|U-\mathbb EU|>t)\le2e^{-t^2/(2v_S)}$ from $t=R/2$;
the polynomial factor is absorbed by doubling the denominator in the
exponent. Put $v=N+rm/4$. Since zero truncation removes only tail values,
\eqref{eq:maskl2} implies
\[
\mathbb E_{z,x,a}(P_z-P_*)^2
 \le64(q-1)v\exp\bigl(-R^2/(16v)\bigr).
\]
If $q=N$, the untruncated score is $\pm2N$, so test $-P_z\ge0$.
If $q<N$, its attainable values are multiples of $4q$; testing
$-P_z\ge2q$ reproduces its negative-sign vote, with a gap of $2q$ on
either side. In both cases disagreement with the vote requires
$|P_z-P_*|\ge2q$. Markov's inequality gives expected disagreement at most
\[
\frac{16v}{q}\exp\bigl(-R^2/(16v)\bigr).
\]
Taking
\[
R=\left\lceil4\sqrt{v\log\frac{16v}{q\delta}}\right\rceil
\]
makes this at most $\delta$. Averaging fixes one mask collection with that
uniform-input error. The resulting circuit is deterministic; masks are
neither supplied inputs nor chosen separately for different inputs.

\emph{The two parameter choices.}
For Part~1 take $q=N$. At balance $v=O(N)$, so
$R=O(\sqrt{N\log(1/\delta)})$. For Part~2 take the least power of two
$q\ge32\varepsilon^2N$ and set $\delta=\varepsilon$.
Then $q<64\varepsilon^2N$, $q\le N$, and the vote's advantage is at least
$\sqrt{q/(8N)}\ge2\varepsilon$. Its truncated circuit has the required
advantage, with $\log(16v/(q\delta))=O(\log(1/\varepsilon))$.
Every bottom has at most $n$ input wires, proving the stated wire bound.
\end{proof}

For each fixed error $0<\delta<1/2$, the upper bound is $O(n^{3/2})$ gates;
the lower bound
of Section~\ref{sec:lower} leaves a factor $\log^2n$. For polynomially small
error or advantage, the new upper bounds have the same asymptotic order as
the earlier union-bound construction. Requiring error below $2^{-n}$ makes
the radius linear in $n$ at balance, so this argument gives no exact
subquadratic circuit.


%===================================================================
\sect{Lower bounds}\label{sec:lower}

\begin{theorem}\label{thm:lower}
Let $C$ be a depth-two threshold circuit of size $s$ with $w$ wires that
$\varepsilon$-agrees with Andreev's function, where
$\varepsilon\ge3(k+2)/\sqrt N$. Then
\[
m\,N\,\varepsilon^2\le 36\,\bigl(k+2+\ln(s+1)\bigr)\cdot s\cdot\sqrt n
\quad\text{and}\quad
m^{3/2}N\varepsilon^2\le 36\,\bigl(k+2+\ln(w+1)\bigr)\cdot w .
\]
If moreover $m\le N$, the first factor can be replaced by $4(k+2)$, giving
$mN\varepsilon^2\le144(k+2)\,s\sqrt n$ and
$m^{3/2}N\varepsilon^2\le144(k+2)\,w$. In the balanced case,
$s=\Omega(\varepsilon^2n^{3/2}/\log^2n)$ and
$w=\Omega(\varepsilon^2n^{5/2}/\log^{5/2}n)$. Taking $\varepsilon=\tfrac12$,
allowed when $N\ge36(k+2)^2$, gives the worst-case bounds.
\end{theorem}

\smalltitle{Plan.}
We follow Kane and Williams~\cite{KW}. Fix all but one bit per block at
random. The restricted circuit computes a lookup in a \emph{random} table
indexed by the $k$ remaining bits, and only a few bottom gates still depend
on those bits (Lemma~\ref{lem:surv}). There are few such restricted circuits
(Lemma~\ref{lem:count}), and none of them correlates well with a random
table (Lemma~\ref{lem:avg}). Averaging over all restrictions proves the
theorem.

\begin{keybox}
\textbf{Notation for Sections~\ref{sec:lower} and~\ref{sec:depth3}.}
$C$ is a depth-two threshold circuit of size $s$ with $w$ wires.
A \emph{restriction} consists of a choice $\sigma$ of one \emph{live}
position in every block, and a \emph{fixing}: values for all other input
bits, table included. For a uniformly random restriction, every fixed bit
is uniform and independent. Write $u\in\bits^k$ for the live bits and
$o_i$ for the parity of the $m-1$ fixed bits of block $i$ (the
\emph{offset}). Then the address is $b=u\oplus o$, and the restricted
function is
\[
u\mapsto x_{u\oplus o}.
\]
A bottom gate \emph{survives} if it still depends on $u$. Let $U$ be the
number of surviving gates.
\end{keybox}

\begin{lemma}[Few survivors]\label{lem:surv}
For a uniformly random restriction, $\mathbb E[U]\le4s\sqrt n/m$. The
expected number of bottom gates that depend on at least two live bits is at
most $4w/m^{3/2}$. Both bounds hold for every fixed table, since they use
only the randomness of the fixed address bits.
\end{lemma}
These are the random restriction lemmas of Kane and
Williams~\cite[Lemmas~1.1 and~3.1]{KW}, with explicit constants.
Appendix~\ref{app:surv} proves both bounds; the proof includes small repairs
to the argument of~\cite{KW} (ties among weights, zero weights, and a strict
versus non-strict inequality). An alternative proof of the first lemma, up to
constants, is in~\cite[App.~A.3]{CTW}; \cite[Lemma~6.1]{KKL}
and~\cite[Lemma~9.1]{CSS} prove analogues with polylogarithmic and $n^{o(1)}$
losses.

\begin{fact}[Roychowdhury, Siu, Orlitsky~\cite{RSO}; see {\cite[Thm.~2.3]{KW}}]\label{fact:rso}
For any fixed functions $f_1,\dots,f_q$ on a finite set $\Omega$, at most
$(|\Omega|+1)^{q+1}$ functions on $\Omega$ have the form
$\ind{\sum_j\omega_jf_j\ge\tau}$.
\end{fact}

\begin{lemma}[Counting restricted circuits]\label{lem:count}
Fix $C$ and $\sigma$, and let $\cH_q$ be the set of functions
$u\mapsto C(\text{restricted})$ over all fixings with at most $q$
survivors. Then
\[
|\cH_q|\le\bigl((2^k+1)(s+1)\bigr)^q\,(2^k+1)^{q+k+1}.
\]
\end{lemma}
\begin{proof}
Non-surviving bottom gates are constants, and are absorbed into the top
threshold. A surviving gate is $\ind{\omega\cdot u\ge\tau}$, where $\omega$ is
its original weight vector restricted to the live positions; only $\tau$
depends on the fixing. For a fixed $\omega$, these functions are nested as
$\tau$ grows, so there are at most $2^k+1$ of them. Pad the list of survivors
with constant functions to length exactly $q$. Then each entry of the list is
one of at most $(2^k+1)(s+1)$ functions: choose one of the at most $s$ gates
or ``constant'', then the threshold. The top gate is a threshold function of
the $q$ entries and the $k$ live bits. By Fact~\ref{fact:rso} with
$|\Omega|=2^k$, there are at most $(2^k+1)^{q+k+1}$ such functions for each
list.
\end{proof}
Kane and Williams~\cite{KW}, and also~\cite{KKL,CSS}, count all threshold
functions of the live bits at this step, which costs about $k^2$ bits per
survivor instead of $O(k+\log s)$.

\begin{lemma}[Agreement bound for all survivor counts at once]\label{lem:avg}
Let $\alpha=\sqrt{N(k+2+\ln(s+1))}$ and $\beta=(k+2)\sqrt N$. Then $C$
agrees with Andreev's function on at most a
\[
\tfrac12+\frac{\alpha\sqrt{\mathbb E[U]}+\beta}{N}+2^{-(k+6)}
\]
fraction of inputs.
\end{lemma}
\begin{proof}
Fix $\sigma$ and the fixed address bits; the table remains uniformly random.
For one function $f$ of $u$ and $\Delta\ge0$, Hoeffding's inequality gives
\[
\Pr\bigl[f\text{ agrees with }u\mapsto x_{u\oplus o}\text{ on}\ge N/2+\Delta\text{ points}\bigr]
\le e^{-2\Delta^2/N}.
\]
Call the fixing \emph{bad} if some $q<N$ and some $f\in\cH_q$ violate this
with $\Delta_q=\alpha\sqrt q+\beta$. Using $\Delta_q^2\ge\alpha^2q+\beta^2$,
$\ln(2^k+1)\le k+1$ and Lemma~\ref{lem:count},
\begin{align*}
\Pr[\text{bad}]&\le\sum_{q\ge0}|\cH_q|\,e^{-2\Delta_q^2/N}
\le\sum_{q\ge0}e^{(2q+k+1)(k+1)+q\ln(s+1)-2q(k+2+\ln(s+1))-2(k+2)^2}\\
&=\sum_{q\ge0}e^{-2q-q\ln(s+1)+(k+1)^2-2(k+2)^2}
\le\sum_{q\ge0}e^{-2q-(k^2+6k+7)},
\end{align*}
which is at most $1.2\,e^{-(k^2+6k+7)}\le2^{-(k+6)}$.

On a good fixing with $U<N$ survivors, the agreement on the $N$ points $u$
is at most $N/2+\alpha\sqrt U+\beta$. This also holds for $U\ge N$, because
then $\alpha\sqrt U\ge\sqrt{2N}\sqrt N>N/2$. Averaging over all restrictions
and using $\mathbb E\sqrt U\le\sqrt{\mathbb E U}$ (Cauchy--Schwarz) proves
the lemma, since the average agreement over all restrictions equals the
agreement over all inputs.
\end{proof}
Kane and Williams~\cite{KW} (and similarly~\cite{KKL,CSS}) instead bound $U$
by Markov's inequality and then use a single survivor count. That loses a
factor of $\varepsilon$; the uniform bound with Cauchy--Schwarz does not.

\begin{proof}[Proof of Theorem~\ref{thm:lower}]
By hypothesis, $\beta/N=(k+2)/\sqrt N\le\varepsilon/3$ and
$2^{-(k+6)}\le\varepsilon/192$. So Lemma~\ref{lem:avg} gives
$\varepsilon/3\le\alpha\sqrt{\mathbb E[U]}/N$. Squaring and inserting
Lemma~\ref{lem:surv},
\[
\varepsilon^2N^2\le9\alpha^2\,\mathbb E[U]\le
9N\bigl(k+2+\ln(s+1)\bigr)\cdot\frac{4s\sqrt n}{m},
\]
which is the first bound.

For wires, count as survivors only the gates that depend on at least two
live bits. A gate depending on at most one live bit computes that bit, its
negation, or a constant, so it can be absorbed into the top gate, which reads
the live bits anyway. Every remaining survivor has at least two input wires,
so there are at most $w$ candidate gates, and Lemmas~\ref{lem:count}
and~\ref{lem:avg} hold with $s$ replaced by $w$. Now use the second bound of
Lemma~\ref{lem:surv}.

Finally, let $m\le N$. If $\ln(s+1)\le3(k+2)$, then $k+2+\ln(s+1)\le4(k+2)$.
Otherwise $s\ge e^{3(k+2)}-1>N^{5/2}\ge mN$, so the claim with $4(k+2)$ holds
trivially. The same argument works for $w$, since $N^{5/2}\ge m^{3/2}N$.
\end{proof}

\begin{remark}[Limits of the method; informal]
After a restriction with one live bit per block, the function
$u\mapsto x_{u\oplus o}$ is an OR of at most $N$ ANDs of literals, which is a
depth-two circuit with $N$ bottom gates. Gates such as $\ind{t_S\ge\theta}$
survive with probability $\Theta(k/\sqrt N)$. So no argument that counts
survivors of a single restriction can give more than about $n^{3/2}/\log n$,
one logarithm above Theorem~\ref{thm:lower}. Kane and Williams already observe
that their restriction lemma is tight for majority gates~\cite[Section~3]{KW}.
Rectangle methods (discrepancy, sign rank) cannot reach $n^{2-o(1)}$ either:
for every split of the inputs between two players, Andreev's function has a
protocol with $O(k)$ bits of communication. One player sends its partial
block parities, and the other then knows $b$ and either announces $x_b$ or
sends $b$ back. Sign rank is also known to be too weak against
$\THR\circ\THR$ in general~\cite{CM}.
\end{remark}

%===================================================================
\sect{Depth three}\label{sec:depth3}

\subsect{Top fan-in for Andreev's function}

A family of depth-two circuits $C_1,\dots,C_M$ is an
\emph{$\varepsilon$-approximate majority} for $A$ if, on every input, at
least a $\tfrac12+\varepsilon$ fraction of the $C_i$ output $A$.

\begin{corollary}\label{cor:fanin}
Let $C_1,\dots,C_M$ be depth-two threshold circuits of size at most $s$
each, and assume $m\le N$.
\begin{enumerate}
\item If they form an $\varepsilon$-approximate majority for $A$ with
      $\varepsilon\ge3(k+2)/\sqrt N$, then
      $mN\varepsilon^2\le144(k+2)\,s\sqrt n$.
\item If on every input strictly more than half of them output $A$, then
      either $M>\sqrt N/(6(k+2))$ or $mN\le576(k+2)\,M^2s\sqrt n$.
\end{enumerate}
In the balanced case, Part~2 says that a majority of $M$ depth-two circuits
of size $s$ computing $A$ needs
$M\ge\min\bigl\{\sqrt N/(6(k+2)),\ \tilde\Omega(n^{3/4}/\sqrt s)\bigr\}$.
The same holds for wires, with $\sqrt n$ omitted and $m$ replaced by
$m^{3/2}$.
\end{corollary}
\begin{proof}
Count the pairs (input, $i$) for which $C_i$ is correct. Every input
contributes at least $(\tfrac12+\varepsilon)M$ pairs, so some $C_i$ is
correct on at least a $\tfrac12+\varepsilon$ fraction of inputs. Apply
Theorem~\ref{thm:lower} (the version for $m\le N$) to that $C_i$. For
Part~2, strictly more than half correct means at least
$\tfrac12+\tfrac1{2M}$, so take $\varepsilon=1/(2M)$; the condition on $M$
is the hypothesis of Theorem~\ref{thm:lower}.
\end{proof}
Kane and Williams state the example ``$n^{1.1}$ gates require top fan-in
$n^{2/15}$'', which corresponds to $M\ge\tilde\Omega(n^{1/2}/s^{1/3})$
(implicit in~\cite[Cor.~1.1]{KW}). With Corollary~\ref{cor:fanin}, a majority
of depth-two circuits of size $n^{1.1}$ that computes $A$ needs top fan-in
$n^{1/5-o(1)}$. Part~2 assumes nothing about how the top majority is built;
it only uses that more than half of the subcircuits are correct.

\subsect{The function \texorpdfstring{$B_n$}{B\_n} (proof outline; formalized in Lean)}

Andreev's function has depth-three majority circuits of size
$O(n)$~\cite[Thm.~1.2]{KW}, so depth-three lower bounds need a harder
function. Kane and Williams~\cite[Section~6]{KW} use the following one.

\begin{keybox}
\textbf{Notation for this subsection.} $k$ and $m$ keep their roles as the
number of address blocks and the block length; $u$, $o$, $\sigma$ and $U$
are as in Section~\ref{sec:lower}, and the address is again $b=u\oplus o$.
The table is replaced by a \emph{seed} $y\in\bits^d$ and a list
$r_z\in\bits^d$, indexed by $z\in\bits^k$, that is \emph{$\eta$-biased}: for
every nonzero $e\in\bits^d$, $\bigl|\sum_z(-1)^{r_z\cdot e}\bigr|\le\eta2^k$.
The function is
\[
B(y,a)=r_b\cdot y\bmod2 ,
\]
and the input length is $n=d+km$. For $f:\bits^k\to\bits$ and
$v\in\{\pm1\}^{2^k}$, the \emph{correlation} is
$\mathrm{corr}(f,v)=2^{-k}\bigl|\sum_z(-1)^{f(z)}v_z\bigr|$; if $f$ agrees with
the sign pattern $v$ on a $\tfrac12+\delta'$ fraction of points, then
$\mathrm{corr}(f,v)\ge2\delta'$. In this subsection, $C$ and $s$ may refer to
the depth-three circuit and its size; this is stated where it applies.
\end{keybox}

Kane and Williams take $d=km$, $k=16\log_2d$, and an explicit list with
$\eta=O(d^{-7})$ from Alon, Goldreich, H{\aa}stad and Peralta~\cite{AGHP}. We
use the same shape with explicit parameters: $k=16p$ blocks, $d=km\le2^p$ and
$\eta\le2^{-7p}$. The powering construction of~\cite{AGHP} gives such lists,
with bias at most $(d-1)2^{-8p}\le2^{-7p}$, so the function is in $\mathsf P$.

\begin{lemma}[Nearly orthogonal vectors]\label{lem:gram}
Let $v_1,\dots,v_J\in\{\pm1\}^D$ satisfy $|\langle v_i,v_j\rangle|\le\eta D$
for $i\ne j$, and let $\xi\in\{\pm1\}^D$. If $|\langle\xi,v_i\rangle|\ge\gamma D$
for every $i$, then $J(\gamma^2-\eta)\le1$.
\end{lemma}
\begin{proof}
Let $v=\sum_i\pm v_i$, with signs chosen so that each
$\langle\xi,\pm v_i\rangle\ge\gamma D$. Then $\langle\xi,v\rangle\ge J\gamma D$
and $\|v\|^2\le JD+J^2\eta D$. Cauchy--Schwarz gives
$J^2\gamma^2D^2\le D(JD+J^2\eta D)$, which rearranges to the claim.
\end{proof}
This is a version of~\cite[Lemma~6.1]{KW}, and a form of the Johnson
bound~\cite[Ch.~3]{Gur}; the sharper form is $J\le(1-\eta)/(\gamma^2-\eta)$.
Apply it with $D=2^k$ and $v_y=\bigl((-1)^{r_z\cdot y}\bigr)_z$: for seeds
$y\ne y'$, $\langle v_y,v_{y'}\rangle$ is the sum of the character of $y+y'$
over the list, so $|\langle v_y,v_{y'}\rangle|\le\eta D$. Hence, if
$\gamma^2>\eta$, a fixed function of $z$ has correlation at least $\gamma$
with at most $1/(\gamma^2-\eta)$ of the vectors $v_y$.

\begin{fact}[Discriminator; {\cite{HMPST}}, {\cite[Prop.~9.5]{CSS}}]\label{fact:disc}
Let $h_1,\dots,h_M:\bits^k\to\bits$ and let $\kappa$ be an integer. If the
function $\ind{h_1+\dots+h_M\ge\kappa}$ agrees everywhere with a function whose
sign pattern $v$ satisfies $2^{-k}|\sum_zv_z|\le\eta$, then some $h_i$ has
$\mathrm{corr}(h_i,v)\ge(1-(M+1)\eta)/M$.
\end{fact}

\needspace{14\baselineskip}
\begin{theorem}\label{thm:B}
Let $p\ge1$, $k=16p$, $d=km\le2^p$, and let the list be $2^{-7p}$-biased.
Write $K_1=272p^2+54p+3$ and $K_2=272p^2+57p+3$.
\begin{enumerate}
\item Let $C$ consist of depth-two threshold circuits, possibly sharing bottom
      gates, with $s$ gates and $w$ wires (the top gate and its wires are not
      counted), and a top gate that outputs $1$ exactly when at least $\kappa$
      of $M$ selected outputs are $1$. The outputs may repeat and are counted
      with multiplicity, and $M\le d^3/4$. If $C$ computes $B$, then
      \[
      d\,(d-K_1)\le4864\,p^2\,s\sqrt{2d}\qquad\text{and}\qquad
      d^{3/2}(d-K_1)\le19456\,p^{5/2}\,w .
      \]
\item If $d>K_2$, every depth-two threshold circuit with $s$ gates and $w$
      wires agrees with $B$ on at most a
      \[
      \tfrac12+2^{-3p}+\frac{2432\,p^2\,s\sqrt{2d}}{d\,(d-K_2)}
      \quad\text{and at most a}\quad
      \tfrac12+2^{-3p}+\frac{9728\,p^{5/2}\,w}{d^{3/2}(d-K_2)}
      \]
      fraction of inputs.
\end{enumerate}
\end{theorem}
\emph{Reading the bounds.} Take $m=\lfloor2^p/(16p)\rfloor$; for $p\ge9$ this
gives $2^{p-1}<d\le2^p$, so $p=\lceil\log_2d\rceil$. Part~1 then gives
$\Omega(n^{3/2}/\log^2n)$ gates and $\Omega(n^{5/2}/\log^{5/2}n)$ wires, and
Part~2 gives agreement $\tfrac12+O(s\log^2n/n^{3/2})+O(n^{-3})$, that is,
$\varepsilon$-agreement needs $\Omega(\varepsilon n^{3/2}/\log^2n)$ gates when
$\varepsilon\ge2\cdot2^{-3p}$. The bounds are nontrivial only for $p\ge17$; at
$p=17$, where $d=130832$, Part~1 certifies only about $10$ gates. The
hypothesis $M\le d^3/4$ costs nothing for majority circuits in the sense of
Kane and Williams: if the subcircuits are distinct then $M\le s$, and when
$s>d^3/4$ Part~1 holds trivially. Chen, Santhanam and Srinivasan use the same
kind of fan-in restriction~\cite[Lemma~9.8]{CSS}.

\begin{proof}[Proof outline with constants]
Let $\gamma=2^{-3p}$. Kane and Williams~\cite[Cor.~6.2]{KW}, and
also~\cite{CSS,KKL}, fix one seed that is bad for all small circuits. We fix the
circuit first, which allows fixed-weight counting. We may assume $s<2^{4p}$
(respectively $w<2^{4p}$), since otherwise the bounds hold trivially; then
$\log_2(s+1)\le4p$. Note that $m=d/(16p)$.
\begin{enumerate}
\item \emph{Survivor budget.} Let $q_0$ be the largest integer $Q$ with
      $(36p+2)Q+K\le d$, where $K=K_1$ in Part~1 and $K=K_2$ in Part~2. Then
      $q_0+1>(d-K)/(36p+2)$.
\item \emph{Few bad seeds.} For each $\sigma$ and each offset $o$, take the
      restricted functions with at most $q_0$ survivors, shifted by $o$. Call a
      seed \emph{bad} if one of them has correlation at least $\gamma$ with
      $v_y$, and also call the zero seed bad. By Lemma~\ref{lem:count},
      $\log_2|\cH_q|\le q(36p+2)+(16p+1)^2$. There are $m^k\le2^{16p^2}$
      choices of $\sigma$ and $2^k$ offsets. By Lemma~\ref{lem:gram}, each
      function correlates with at most $1/(\gamma^2-\eta)\le2^{6p+1}$ seeds.
      Together,
      \[
      \Pr[\text{bad seed}]\le2^{(36p+2)q_0+272p^2+54p+2-d}.
      \]
      With $K=K_1$ this is at most $\tfrac12$; with $K=K_2$ it is at most
      $2^{-3p-1}=\gamma/2$.
\item \emph{Few survivors.} Lemma~\ref{lem:surv} gives
      $\mathbb E[U]\le4s\sqrt{2d}/m=64p\,s\sqrt{2d}/d$, the expectation being
      over the restriction and the seed together. By Markov's inequality and
      step~1,
      \[
      \Pr[U>q_0]\le\frac{64p(36p+2)\,s\sqrt{2d}}{d(d-K)}
      \le\frac{2432\,p^2\,s\sqrt{2d}}{d(d-K)} .
      \]
      For wires, $\mathbb E[U]\le4w/m^{3/2}=256p^{3/2}w/d^{3/2}$, and
      $256p^{3/2}(36p+2)\le9728p^{5/2}$.
\item \emph{Part 2.} On a good seed with at most $q_0$ survivors, the circuit
      agrees with $B$ on at most a $\tfrac12(1+\gamma)$ fraction of the
      restricted inputs. So the agreement is at most
      $\tfrac12(1+\gamma)+\Pr[\text{bad seed}]+\Pr[U>q_0]\le\tfrac12+\gamma+\Pr[U>q_0]$.
\item \emph{Part 1.} Suppose $\Pr[U>q_0]<\tfrac12$. Together with step~2, some
      restriction has a good seed and at most $q_0$ survivors. On it, $C$
      computes exactly the function of $u$ whose sign pattern is
      $u\mapsto(v_y)_{u\oplus o}$. A good seed is nonzero, so this target has
      bias at most $\eta$. By
      Fact~\ref{fact:disc}, some subcircuit has correlation at least
      $(1-(M+1)\eta)/M>\gamma$ with it, since $M\le d^3/4\le2^{3p}/4$. That
      contradicts the choice of the seed. Hence $\Pr[U>q_0]\ge\tfrac12$, which
      by step~3 gives the factor $4864=2\cdot2432$, and $19456=2\cdot9728$ for
      wires.\qedhere
\end{enumerate}
\end{proof}
The formal proof (Section~\ref{sec:lean}) follows these steps.

Part~2 is not essentially new: Kane and Williams' proof
of~\cite[Cor.~6.2]{KW} combined with Markov's inequality gives the same form
with an extra logarithm, although they do not state it. They do say that they
cannot prove correlation bounds of the form $n^{-\omega(1)}$ for $B_n$, and
that their approach cannot go below correlation $1/\sqrt n$ for circuits with
about $n$ gates~\cite[Section~6]{KW}. Part~2 agrees: for $s\approx n$ it
gives about $\log^2n/\sqrt n$. What is new is the saved logarithm, obtained
by fixing the circuit first.

%===================================================================
\sect{Formalization status}\label{sec:lean}

The Lean~4 proofs are in Complexitylib,
\url{https://github.com/cslib-community/complexitylib}. Pull requests
\#2377, \#2813 and \#2814 are merged; \#2816 remains open as of
7 November 2026. The refined exact and mean-square upper bounds proved in
this note are not yet formalized. The Lean statements are explicit finite
inequalities; their asymptotic corollaries are derived in the text.

\begin{center}
\small
\renewcommand{\arraystretch}{1.2}
\begin{tabular}{@{}P{0.30\textwidth}P{0.23\textwidth}P{0.40\textwidth}@{}}
\toprule
Result & Status & Where\\
\midrule
Proposition~\ref{prop:basic} & formalized, merged & PR \#2813; \texttt{Andreev.lean}\\
Theorem~\ref{thm:exact} & proof here; not formalized & refined exact score\\
Theorem~\ref{thm:avgupper} & proof here; not formalized & mean-square mask argument\\
Earlier windowed upper bounds & formalized, merged & \texttt{AndreevWindowed.lean}, \texttt{AndreevAverage.lean}\\
Lemma~\ref{lem:surv} & formalized, merged & \texttt{RandomRestriction.lean}, \texttt{ThresholdCounting.lean}\\
Lemmas~\ref{lem:count},~\ref{lem:avg}, Theorem~\ref{thm:lower} & formalized, merged & PR \#2377; \texttt{AndreevLowerBound.lean}\\
Corollary~\ref{cor:fanin} & formalized, merged & PR \#2814; \texttt{AndreevApproxMajority.lean}\\
Lemma~\ref{lem:gram}, Theorem~\ref{thm:B} & formalized, open PR & PR \#2816; \texttt{KaneWilliamsB.lean}, \texttt{NearlyOrthogonal.lean}\\
\bottomrule
\end{tabular}
\end{center}

The merged average-case code proves the earlier union-bound window
construction, using $\varphi(X+t)+\varphi(X-t)=2X(-1)^t$. It does not prove
the sharper logarithmic factors in Theorem~\ref{thm:avgupper}. Likewise the
merged exact upper bound is Proposition~\ref{prop:basic}, not
Theorem~\ref{thm:exact}. The sign-coherent appendix is not formalized.

For Corollary~\ref{cor:fanin}, the top majority is modelled as a hypothesis
on the number of correct subcircuits. For Theorem~\ref{thm:B}, the seed
length $d$ is called $n$ and the block length $m$ is called $\ell$ in Lean;
the circuit has $M$ outputs, followed by the gate ``at least $\kappa$ of the
$M$ outputs''. Part~2 uses $K_2$ and Part~1 uses $K_1$ as in the text.

\smalltitle{Acknowledgements.}
The results, the formalization and this note were developed with AI
assistance.

\appendix
\renewcommand{\thesection}{\Alph{section}}
\sect{Proof of Lemma~\ref{lem:surv}}\label{app:surv}

\begin{fact}[Erd\H{o}s~\cite{Erdos}]\label{fact:lo}
Let $z_1,\dots,z_R$ be independent uniform bits, let $\omega_1,\dots,\omega_R$
satisfy $|\omega_j|\ge\mu>0$, and let $I$ be a half-open interval of length
$\mu$. Then
$\Pr\bigl[\sum_j\omega_jz_j\in I\bigr]\le\binom{R}{\lfloor R/2\rfloor}2^{-R}\le1/\sqrt R$.
\end{fact}

Fix one bottom gate $\ind{\sum_j\omega_jz_j\ge\tau}$, where $j$ runs over
all input positions. Let $\Lambda$ be the set of address positions with
$\omega_j\ne0$, so $|\Lambda|$ is at most the number of the gate's input wires.
For $j\in \Lambda$ let $\rho(j)$ be the number of $j'\in \Lambda$ with
$|\omega_{j'}|\ge|\omega_j|$, including $j$ itself. Fix the table bits; they
only shift the threshold.

\begin{lemma}\label{lem:onerestr}
Fix the live positions. Over the random fixed address bits, the gate is
non-constant on the live bits with probability at most
$2\sum_{j}\rho(j)^{-1/2}$, the sum over the live positions $j\in \Lambda$.
\end{lemma}
\begin{proof}
Let $\Phi$ be the contribution of the fixed bits. The restricted gate is
$\ind{\sum_{j\text{ live}}\omega_ju_j\ge\tau-\Phi}$, and it is non-constant
exactly when $\tau-\Phi$ lies in a half-open interval of length
$\sum_{j\text{ live}}|\omega_j|$. Split this interval into consecutive
half-open pieces of lengths $|\omega_j|$, one for each live $j\in \Lambda$.

For a live $j\in \Lambda$, let $R_j$ be the number of \emph{fixed} positions
$j'\in \Lambda$ with $|\omega_{j'}|\ge|\omega_j|$. If $R_j\ge\rho(j)/4$ for every live
$j\in \Lambda$, then Fact~\ref{fact:lo}, applied to these $R_j$ positions and
conditioned on all other fixed bits, bounds the probability of piece $j$ by
$R_j^{-1/2}\le2\rho(j)^{-1/2}$, and summing over the pieces proves the lemma.

Otherwise some live $j$ has $R_j<\rho(j)/4$, so more than $3\rho(j)/4$ live
positions $j'\in \Lambda$ have $|\omega_{j'}|\ge|\omega_j|$, and each of them has
$\rho(j')\le\rho(j)$. Then the claimed bound is at least
$\tfrac34\rho(j)\cdot2\rho(j)^{-1/2}\ge\tfrac32$, which exceeds any
probability.
\end{proof}

\begin{proof}[Proof of Lemma~\ref{lem:surv}]
Each position of a block is live with probability $1/m$. Sorting the values
$\rho(j)$, the $t$-th smallest is at least $t$. So the gate survives with
probability at most
\[
\frac2m\sum_{j\in \Lambda}\rho(j)^{-1/2}\le\frac2m\sum_{t=1}^{|\Lambda|}t^{-1/2}
\le\frac{4\sqrt{|\Lambda|}}{m}\le\frac{4\sqrt n}{m}.
\]
Summing over the at most $s$ bottom gates gives the first bound.

A gate that depends on at least two live bits is non-constant and has at
least two live positions in $\Lambda$. Positions in different blocks are live
independently, so for $j\in \Lambda$,
\[
\Pr[\,j\text{ and another position of }\Lambda\text{ are live}\,]
\le\min\Bigl\{\frac1m,\frac{|\Lambda|}{m^2}\Bigr\}\le\frac{\sqrt{|\Lambda|}}{m^{3/2}}.
\]
By Lemma~\ref{lem:onerestr}, the gate depends on at least two live bits with
probability at most
\[
2\sum_{j\in \Lambda}\frac{\sqrt{|\Lambda|}}{m^{3/2}}\,\rho(j)^{-1/2}
\le\frac{2\sqrt{|\Lambda|}}{m^{3/2}}\cdot2\sqrt{|\Lambda|}=\frac{4|\Lambda|}{m^{3/2}}.
\]
Summing over the bottom gates, and using that $\sum|\Lambda|$ is at most the number
of wires $w$, gives the second bound.
\end{proof}

The formal proofs are the Lean theorems
\texttt{card\_not\_forcedToConst\_ltfEval\_le} (Lemma~\ref{lem:onerestr}, there
with all non-live positions random),
\texttt{sum\_circuitUnforcedConstCount\_blockLiveSet\_le} and
\texttt{sum\_circuitUnforcedAtMostOneCount\_blockLiveSet\_le}.

\sect{Note: near-quadratic lower bounds for sign-coherent circuits}\label{app:coherent}

\begin{keybox}
\textbf{Status.} This appendix records a recent result together with
everything needed to reproduce it. The proofs below are complete and have
been checked by an independent referee and by computer on small cases. The
constants are not optimised, and the result is not yet formalized in Lean.
We plan to tighten and integrate it later.
\end{keybox}

\smalltitle{The class.}
A bottom gate is \emph{sign-coherent} if, inside every address block, its
weights are either all $\ge0$ or all $\le0$. Its table weights and the top
gate are arbitrary. Every gate that sees a block only through $|a_i|$ is
sign-coherent, and so is every gate that is monotone in the address bits.
All gates of the circuit in Theorem~\ref{thm:exact} are of this kind.

\begin{theorem}\label{thm:coherent}
Let $C$ be a depth-two threshold circuit for $A$ whose $S$ bottom gates are
all sign-coherent. Put $D_0=\max\{0,N-k-1-\log_2(S+1)\}/\bigl(k+\log_2(S+1)\bigr)$. Then
\[
S\ \ge\ \frac{m^k}{m^k-(m-1)^k}\,D_0\ \ge\ \frac mk\,D_0 .
\]
Consequently $S\ge\min\{N^2,\ m(N-3k-1)/(3k^2)\}$, and in the balanced case
\[
S\ \ge\ \frac{N(N-3k-1)}{3k^3}\ =\ \Omega\!\left(\frac{n^2}{\log^3 n}\right).
\]
\end{theorem}

So within this class the circuit of Theorem~\ref{thm:exact} is optimal up to a
polylogarithmic factor. A polynomial-exponent improvement below quadratic therefore requires
gates whose weights have mixed signs inside a block.

\smalltitle{Plan.}
We walk through all unit cubes of the grid of block weights. On each cube,
the circuit reveals the whole table, so many gates must be non-constant there
(Lemma~\ref{lem:entropy}). A sign-coherent gate is non-constant on only a
$k/m$ fraction of the cubes (Lemma~\ref{lem:diag}). Comparing the two counts
gives the theorem.

\smalltitle{The cube restriction.}
Fix an order of the positions inside every block. For a vector
$s\in\{0,\dots,m\}^k$ let $a(s)$ be the address part in which block $i$ has
ones exactly at its first $s_i$ positions. A \emph{cube} is a vector
$c\in\{0,\dots,m-1\}^k$; there are $m^k$ cubes. For $z\in\bits^k$ the input
$(x,a(c+z))$ differs from $(x,a(c))$ exactly in the live positions
$c_i+1$ of the blocks $i$ with $z_i=1$, and its address is
$(c\bmod2)\oplus z$. Hence
\[
z\mapsto A\bigl(x,a(c+z)\bigr)=x_{(c\bmod 2)\oplus z}
\]
determines the whole table $x$. A bottom gate \emph{splits} the cube $c$ for
the table $x$ if it is non-constant in $z$ on the inputs $(x,a(c+z))$. Let
$d_c(x)$ be the number of bottom gates that split $c$ for $x$.

\begin{lemma}[Many gates split each cube]\label{lem:entropy}
Let $C$ be any depth-two threshold circuit for $A$ with $S$ bottom gates, and
let $x$ be uniformly random. Then $\mathbb E_x[d_c(x)]\ge D_0$ for every cube
$c$.
\end{lemma}
\begin{proof}
Fix $c$. On the inputs $(x,a(c+z))$, bottom gate $j$ computes
$\ind{\lambda_j\cdot z\ge\tau_j}$, where $\lambda_j\in\mathbb R^k$ lists its
weights on the live positions and only the number $\tau_j$ depends on $x$.
Likewise the top gate computes
$\ind{\sum_jv_jg_j+\mu\cdot z\ge\kappa}$, where $v_j$ are its weights on the
bottom gates and $\mu$ its weights on the live positions; the gates that
do not split $c$ are constant and are absorbed, together with the table and
the fixed address bits, into $\kappa$. So only $\tau_j$ and $\kappa$ depend
on $x$.

Therefore the function $z\mapsto x_{(c\bmod2)\oplus z}$, and with it $x$, is
determined by: the set $G$ of splitting gates; for each $j\in G$ the
non-constant function $\ind{\lambda_j\cdot z\ge\tau_j}$, one of at most $2^k-1$
(only the position of $\tau_j$ among the $\le2^k$ values $\lambda_j\cdot z$
matters); and the position of $\kappa$, one of at most $2^k+1$. With
$D=|G|=d_c(x)$, entropy gives
\[
N=H(x)\le H(D)+\mathbb E\log_2\Bigl[\tbinom SD(2^k-1)^D(2^k+1)\Bigr]
\le\log_2(S+1)+\mathbb E[D]\,(\log_2(S+1)+k)+k+1,
\]
which rearranges to the claim; a negative numerator is replaced by the
trivial lower bound zero.
\end{proof}

\begin{lemma}[A coherent gate splits few cubes]\label{lem:diag}
For every table $x$, a sign-coherent gate splits at most $m^k-(m-1)^k$ cubes.
\end{lemma}
\begin{proof}
Let $\varepsilon_i=+1$ if the gate's weights in block $i$ are $\ge0$ and
$\varepsilon_i=-1$ otherwise. Its linear form on $(x,a(s))$ is a constant plus
$\sum_iP_i(s_i)$, where $P_i(s_i)$ is the sum of its weights on the first
$s_i$ positions of block $i$. So the form $F(s)$ does not decrease when
$s_i$ moves by one step in direction $\varepsilon_i$.

On the cube $c$ the form is smallest at $c^-=c+(\mathbf1-\varepsilon)/2$ and
largest at $c^+=c+(\mathbf1+\varepsilon)/2$. The gate with threshold $\theta$
splits $c$ exactly when $F(c^-)<\theta\le F(c^+)$.

Group the cubes into \emph{lines} $c,\,c+\varepsilon,\,c+2\varepsilon,\dots$.
For consecutive cubes on a line, the largest vertex of one is the smallest
vertex of the next, so the intervals $(F(c^-),F(c^+)]$ along a line are
disjoint, and $\theta$ lies in at most one of them. Each line starts at a cube
$c$ for which $c-\varepsilon$ is not a cube, and exactly $(m-1)^k$ cubes $c$
have $c-\varepsilon$ a cube. So there are $m^k-(m-1)^k$ lines.
\end{proof}

\begin{proof}[Proof of Theorem~\ref{thm:coherent}]
Sum Lemma~\ref{lem:entropy} over all cubes and compare with
Lemma~\ref{lem:diag}: $m^kD_0\le\sum_c\mathbb E[d_c]\le S\,(m^k-(m-1)^k)$.
Then use $1-(1-1/m)^k\le k/m$. If $S<N^2$ then $\log_2(S+1)\le2k$ and
$D_0\ge(N-3k-1)/(3k)$. In the balanced case $m=N/k$; this requires $k$ to be
a power of two.
\end{proof}

\begin{remark}
(1) For gates that see each block only through $|a_i|$, a slab argument with
Ball's cube-slicing bound improves the denominator to $6\sqrt2\,k^{5/2}$ for
$k\ge8$. (2) Lemma~\ref{lem:diag} is tight for all $k,m$. The fact that a
unate function changes at most once along a chain compatible with its
orientation is classical (Impagliazzo, Paturi and Saks~\cite{IPS},
Lemma~2.1); the cube restriction and the quadratic bound appear to be new,
but our literature search was limited. (3) Theorem~\ref{thm:coherent} is
asymptotic. It exceeds $n^{3/2}$ only from about $k=37$, and the method
(Lemma~\ref{lem:entropy} plus a uniform bound as in Lemma~\ref{lem:diag})
cannot give more than $O(n^2/\log n)$. (4) Lemma~\ref{lem:entropy} holds for
all circuits. For general gates, the survival-type arguments of
Section~\ref{sec:lower} cannot exceed $O(n^{3/2})$: under every
distribution of restrictions that leaves one live bit per block, a gate with
random $\pm1$ weights is non-constant with probability at least
$1/(2\sqrt{2km})$. Gates with mixed signs inside a block are exactly where
the gap between $n^{3/2}$ and $n^2$ lies.
\end{remark}

\smalltitle{Two address blocks.}
For $k=2$ the bounds are linear in $m$ on both sides.

\begin{proposition}\label{prop:k2}
Let $k=2$. (a) Every circuit for $A$ with sign-coherent bottom gates has at
least $m-1$ bottom gates. (b) There is a circuit with $4\lceil m/2\rceil+m$
bottom gates, each seeing the blocks only through $s_i=|a_i|$.
\end{proposition}
\begin{proof}[Proof sketch of (a)]
Take the table with $x_c=c_1\oplus c_2$. On every cube (here a unit square)
the restricted circuit computes XOR or its negation. If every bottom gate were
constant or a literal there, the top gate would compute XOR as a threshold
function, which is impossible. So on every square some gate cuts off exactly
one corner. After reflecting $s_i\mapsto m-s_i$ where needed, the region where
a coherent gate is $1$ is an up-set of the grid $\{0,\dots,m\}^2$. Its
boundary is a monotone lattice path, and a corner is cut exactly where the
path turns. On the $4(m-1)$ squares at the border of the grid such a path
turns at most four times (the \emph{Ring Lemma}). So at least $m-1$ gates are needed.
\end{proof}
\begin{proof}[Construction for (b)]
Write $b=b_1+2b_2$ for the table index and $J=\lceil m/2\rceil$. The bottom
gates are, for $1\le j\le J$,
\[
\ind{s_1+x_0\ge2j},\ \ind{s_1+x_2\ge2j},\ \ind{s_2+x_0\ge2j},\ \ind{s_2+x_1\ge2j}
\]
with top weights $-4,-2,-4,-2$, and, for $1\le j\le m$, $\ind{s_1+s_2+x_0\ge2j}$
with top weight $-4$. The top gate outputs $\ind{F\ge0}$ with
\[
F=\sum_j v_jg_j+5s_1+5s_2+8x_0+2x_1+\tfrac32x_2+x_3-\tfrac{13}2 .
\]
For $0\le t\le T$ we have $t-2\sum_{j\le\lceil T/2\rceil}\ind{t+y\ge2j}=(-1)^y(t\bmod2)$
for every bit $y$, so $F$ depends on $s$ only through $b$. This leaves a
finite check, and we verified $\ind{F\ge0}=x_b$ for all inputs with exact
arithmetic for $m\le40$.
\end{proof}

The earlier general constructions give about $12m$ gates for $k=2$. Exact
solver data for gates that see blocks only through $|a_i|$ give minima $3,3$
for $m=1,2$ and lower bounds $5$ for $m=3,4,5$ and $6$ for $m=6,8,10$.

\smalltitle{Reproducing.}
Lemma~\ref{lem:diag} can be checked by brute force: enumerate all cubes and
compute $F(c^\pm)$ for random sign-coherent weight vectors, including zero
weights and ties, and random orders. Equality is attained by a gate with one
large weight on the last position of each block. For the bound of (a), the
Ring Lemma can be checked over all up-sets of $\{0,\dots,m\}^2$ for
$m\le11$. The construction of (b) is checked as described. All three checks
were run independently twice.

\renewcommand{\refname}{\texorpdfstring{\color{ink}References}{References}}
\begin{thebibliography}{99}\setlength{\itemsep}{0pt}\footnotesize

\bibitem{AGHP} N.~Alon, O.~Goldreich, J.~H{\aa}stad, and R.~Peralta.
\emph{Simple constructions of almost $k$-wise independent random variables}.
Random Structures \& Algorithms 3(3):289--304, 1992.
\bibitem{ACW} J.~Alman, T.~M. Chan, and R.~Williams.
\emph{Polynomial Representations of Threshold Functions and Algorithmic
Applications}. FOCS 2016, pages 467--476.
\bibitem{CM} A.~Chattopadhyay and N.~S. Mande.
\emph{Weights at the Bottom Matter When the Top is Heavy}.
\href{https://eccc.weizmann.ac.il/report/2017/083/}{ECCC TR17-083}, 2017.
\bibitem{CSS} R.~Chen, R.~Santhanam, and S.~Srinivasan.
\emph{Average-Case Lower Bounds and Satisfiability Algorithms for Small
Threshold Circuits}. Theory of Computing 14(9):1--55, 2018.
\bibitem{CTW} L.~Chen, A.~Tal, and Y.~Wang.
\emph{Super-quadratic Lower Bounds for Depth-2 Linear Threshold Circuits}.
STOC 2026. \href{https://eccc.weizmann.ac.il/report/2026/039/}{ECCC TR26-039}.
\bibitem{Erdos} P.~Erd\H{o}s. \emph{On a lemma of Littlewood and Offord}.
Bull.\ Amer.\ Math.\ Soc.\ 51:898--902, 1945.
\bibitem{Gur} V.~Guruswami. \emph{List Decoding of Error-Correcting Codes}.
LNCS 3282, Springer, 2004.
\bibitem{HP} K.~A. Hansen and V.~V. Podolskii. \emph{Polynomial threshold
functions and Boolean threshold circuits}. Information and Computation
240:56--73, 2015.
\bibitem{HMPST} A.~Hajnal, W.~Maass, P.~Pudl\'ak, M.~Szegedy, and
G.~Tur\'an. \emph{Threshold circuits of bounded depth}. J.\ Comput.\ Syst.\
Sci.\ 46(2):129--154, 1993.
\bibitem{IPS} R.~Impagliazzo, R.~Paturi, and M.~Saks. \emph{Size--depth tradeoffs for threshold circuits}. SIAM J.\ Comput.\ 26(3):693--707, 1997.

\bibitem{KKL} V.~Kabanets, D.~M. Kane, and Z.~Lu.
\emph{A Polynomial Restriction Lemma with Applications}. STOC 2017.
\href{https://eccc.weizmann.ac.il/report/2017/026/}{ECCC TR17-026}.
\bibitem{KW} D.~M. Kane and R.~Williams.
\emph{Super-Linear Gate and Super-Quadratic Wire Lower Bounds for Depth-Two
and Depth-Three Threshold Circuits}. STOC 2016.
\href{https://arxiv.org/abs/1511.07860}{arXiv:1511.07860v1}.
\bibitem{Nag} D.~Nag. \emph{Almost-Everywhere Near-Cubic Wire Lower Bounds for
$\mathsf{SYM}\circ\THR$ and $\THR\circ\THR$}. Preprint,
\href{https://eccc.weizmann.ac.il/report/2026/167/}{ECCC TR26-167}, 2026.
\bibitem{RSO} V.~P. Roychowdhury, K.-Y. Siu, and A.~Orlitsky.
\emph{Neural models and spectral methods}. In Theoretical Advances in Neural
Computation and Learning, pages 3--36, 1994.
\bibitem{Tam} S.~Tamaki. \emph{A Satisfiability Algorithm for Depth Two
Circuits with a Sub-Quadratic Number of Symmetric and Threshold Gates}.
\href{https://eccc.weizmann.ac.il/report/2016/100/}{ECCC TR16-100}, 2016.
\end{thebibliography}
\end{document}
